% !TeX root = ../thesis.tex

\chapter{工程背景与设计条件}

\section{项目来源与工艺背景}

江苏斯尔邦石化有限公司甲醇制烯烃及烯烃下游一体化一期项目配套冷媒膨胀罐（设备位号：22-V-1810），隶属于 $100\ \text{kt/a}$ 釜式法 $+$ $200\ \text{kt/a}$ 管式法 LDPE/EVA 联合装置（18 单元）。LDPE/EVA 高压聚合工艺的循环冷却水系统在启停、负荷调节及事故工况下，冷媒温度波动引起介质体积变化，需设置膨胀罐以吸纳体积增量、平抑压力脉动并维持循环泵吸入条件\cite{guocheng2020}。

装置建于江苏省连云港市徐圩新区石化产业基地。连云港地处暖温带季风气候区，极端最高气温 $38.5$ ℃，极端最低气温 $-12.3$ ℃。设备露天布置，承受风载荷、地震载荷及环境温度交变作用，设计使用年限 15 年。

膨胀罐在工艺系统中的功能定位明确为三项：其一，利用罐体上部氮气覆盖层的可压缩性，缓冲循环水系统因温度变化引起的压力波动；其二，在变工况条件下向系统补充或回收冷媒，维持循环回路总容积恒定；其三，通过自由液面实现气液分离，防止气体夹带进入循环泵造成气蚀。

\section{设计依据与标准体系}

本设备按固定式压力容器进行设计、制造、检验与验收，遵循的标准体系如表 \ref{tab:main-standards} 所示。

\begin{table}[!htb]
  \centering
  \caption{设计遵循的标准体系}
  \label{tab:main-standards}
  \small\renewcommand{\arraystretch}{1.1}
  \setlength{\tabcolsep}{7pt}
  \begin{tabularx}{\textwidth}{@{}c >{\centering\arraybackslash}X >{\centering\arraybackslash}X@{}}
    \toprule
    \textbf{序号} & \textbf{标准编号} & \textbf{标准名称} \\
    \midrule
    1  & GB/T 150.1～150.4—2024 & 压力容器（主设计标准） \\
    2  & TSG 21—2016           & 固定式压力容器安全技术监察规程 \\
    3  & GB/T 713.1—2023       & 承压设备用钢板和钢带 第1部分：碳钢和低合金钢 \\
    4  & GB/T 709—2019         & 热轧钢板和钢带的尺寸、外形、重量及允许偏差 \\
    5  & GB/T 25198—2010       & 压力容器封头 \\
    6  & NB/T 47065.1-5—2018   & 容器支座 \\
    7  & NB/T 47020～47027—2012 & 压力容器法兰、垫片、紧固件 \\
    8  & HG/T 20592～20635—2009 & 钢制管法兰、垫片和紧固件 \\
    9  & NB/T 11025—2022       & 补强圈 \\
    10 & HG/T 21514—2014       & 钢制人孔和手孔的类型与技术条件 \\
    11 & NB/T 47013.1～.14—2015 & 承压设备无损检测 \\
    12 & HG/T 20580—2020       & 钢制化工容器设计基础规范 \\
    13 & HG/T 20585—2020       & 钢制低温压力容器技术规范 \\
    14 & NB/T 47041—2014       & 塔式容器 \\
    15 & HG 20660—2000         & 压力容器中化学介质毒性危害和爆炸危险程度分类 \\
    \bottomrule
  \end{tabularx}
\end{table}

原始工艺数据取自中国石化工程建设有限公司编制的设备机械数据表（文件号：3767-CD-SE-0V1810\_IS1）\cite{sierbang-datasheet}。设计参数中内径、容积、压力、温度、腐蚀裕量等直接采用数据表给定值；结构型式、材料牌号、计算方法按 GB/T 150—2024 系列标准选取和确定。

\section{操作介质与容器类别}

\subsection{操作介质与介质组别}

本设备操作介质为冷却水（冷媒），在装置正常运行时工作温度 $155$ ℃，设计温度 $170$ ℃。按 HG 20660—2000 的毒性危害与爆炸危险程度分类\cite{hg20660-2000}，冷却水属于无毒、不易燃介质，划入第二组介质\cite{tsg21-2016}。

\subsection{容器类别划分}

依据 TSG 21—2016 附件 A 的分类准则\cite{tsg21-2016}，固定式压力容器按介质组别和压力-容积乘积划分为 Ⅰ、Ⅱ、Ⅲ 三类：

\begin{equation}
PV = P \times V = 0.80 \times 5.5 = 4.4\ \text{MPa}\cdot\text{m}^3
\label{eq:pv-product}
\end{equation}

对照 TSG 21—2016 附件 A 第二组介质分类图，$PV = 4.4\ \text{MPa}\cdot\text{m}^3 < 50\ \text{MPa}\cdot\text{m}^3$，归属第 Ⅰ 类压力容器。本设备按 Ⅰ 类容器进行设计监管，A、B 类焊接接头允许按局部射线检测执行\cite{tsg21-2016}（检测比例 $\ge 20\%$ 且 $\ge 250$ mm），对应焊接接头系数 $\phi = 0.85$（详见第 2 章）。

\section{操作工况与设计条件}

\subsection{操作工况}

设备在两种操作工况下运行。正常操作工况：工作压力 $0.60$ MPa(g)，工作温度 $155$ ℃；温水充液工况：工作温度 $100$ ℃，用于装置开工及检修后系统充液排气操作。此外，设备须具备全真空运行能力，以应对蒸汽冷凝造成系统负压的极端工况。

\subsection{设计条件的确定}

设计压力和设计温度在操作工况基础上附加安全裕量确定，依照 GB/T 150.1—2024 第 3.1 节的原则执行\cite{gbt150-2024}。

设计压力 $P = 0.80$ MPa(g) 的确定：除覆盖正常操作压力 $0.60$ MPa(g) 外，附加操作波动裕量 $0.20$ MPa 以保证压力安全阀的设定值与操作压力之间有合理的压差带；同时计入全真空外压条件（FV，$P_{\text{ext}} = -0.1013$ MPa(g)），设备壳体须同时满足内压强度和外压稳定性的双重约束。

设计温度 $T = 170$ ℃ 的确定：正常操作最高工作温度 $155$ ℃，附加工艺波动裕量 $15$ ℃（含换热不确定因素及控制系统响应延迟），取设计温度 $170$ ℃。该温度高于温水充液工况温度 $100$ ℃，故以 $170$ ℃ 作为各承压元件强度计算的统一金属设计温度。全真空工况对应金属温度 $100$ ℃（蒸汽冷凝工况），用于外压稳定性校核时弹性模量的取值。

最低设计金属温度（MDMT）$-15$ ℃ 按项目所在地极端最低气温 $-12.3$ ℃ 附加 $-2.7$ ℃ 裕量确定，用于材料低温韧性评定\cite{hgt20585-2020}。本 MDMT 高于 Q345R 正火态的最低允许使用温度 $-20$ ℃，材料低温韧性合格，无需额外进行冲击试验。

\section{设计参数汇总}

综合设备机械数据表和标准规范要求，本冷媒膨胀罐（22-V-1810）的全部设计参数汇总如表 \ref{tab:design-summary}。

\begin{table}[!htb]
  \centering
  \caption{设计参数汇总}
  \label{tab:design-summary}
  \small\renewcommand{\arraystretch}{1.1}
  \setlength{\tabcolsep}{6pt}
  \begin{tabularx}{\textwidth}{@{}c >{\centering\arraybackslash}X >{\centering\arraybackslash}X@{}}
    \toprule
    \textbf{序号} & \textbf{参数名称} & \textbf{数值与说明} \\
    \midrule
    1  & 设备位号 / 名称             & 22-V-1810 / 冷媒膨胀罐 \\
    2  & 压力容器类别                & 第 Ⅰ 类（TSG 21—2016） \\
    3  & 操作介质 / 介质组别          & 冷却水（冷媒）/ 第二组介质，无毒、不易燃 \\
    4  & 工作压力（正常操作）         & $0.60$ MPa(g) \\
    5  & 设计压力                    & $0.80$ MPa(g)，含全真空要求（FV） \\
    6  & 工作温度                    & $155$ ℃（正常操作）/ $100$ ℃（温水充液） \\
    7  & 设计温度                    & $170$ ℃ \\
    8  & 全真空工况金属温度           & $100$ ℃ \\
    9  & 最低设计金属温度（MDMT）     & $-15$ ℃ @ $0.85$ MPa(g) \\
    10 & 罐体内径 $D_i$              & $1800$ mm \\
    11 & 筒体圆筒段长度 $L$           & $1700$ mm（切线至切线，不含封头直边段） \\
    12 & 有效容积 $V$                 & $5.5$ m³ \\
    13 & 最大介质密度                & $1000$ kg/m³（满液校核工况） \\
    14 & 壳体/封头材料               & Q345R 正火钢板（GB/T 713.1—2023） \\
    15 & 接管材料                    & 20 号钢无缝钢管（GB 6479—2013） \\
    16 & 焊接接头系数 $\phi$          & $0.85$（A、B 类焊缝，局部 RT，检测比例 $\ge 20\%$） \\
    17 & 钢板厚度负偏差 $C_1$         & $0.30$ mm（GB/T 709—2019 B类偏差，固定下偏差$-0.30$ mm） \\
    18 & 腐蚀裕量 $C_2$               & $3.0$ mm（壳体/封头）；$2.0$ mm（裙座） \\
    19 & 设计使用年限                & 15 年 \\
    20 & 基本风压 $q_0$               & $450$ N/m²（连云港，50 年一遇，B 类粗糙度） \\
    21 & 抗震设防烈度 / 基本加速度    & 7 度 / $0.10g$ \\
    22 & 设计地震分组 / 场地类别      & 第一组 / Ⅱ 类 \\
    \bottomrule
  \end{tabularx}
\end{table}

\section{技术特性表}

按 HG/T 20580—2020 和 TSG 21—2016 的要求，编制冷媒膨胀罐技术特性表（表 \ref{tab:tech-spec}），作为设计说明书和竣工图的首页依据。

\begin{table}[!htb]
  \centering
  \caption{冷媒膨胀罐 V-1810 技术特性表}
  \label{tab:tech-spec}
  \small\renewcommand{\arraystretch}{1.15}
  \setlength{\tabcolsep}{6pt}
  \begin{tabularx}{\textwidth}{@{}l >{\raggedright\arraybackslash}X l >{\raggedright\arraybackslash}X@{}}
    \toprule
    \multicolumn{4}{c}{\textbf{技术特性表}} \\
    \midrule
    设备位号         & 22-V-1810            & 设备名称     & 冷媒膨胀罐 \\
    容器类别         & 第 Ⅰ 类              & 设计标准     & GB/T 150.1$\sim$150.4—2024 \\
    设计使用年限     & 15 年                 & 监管规程     & TSG 21—2016 \\
    \midrule
    设计压力（内压） & $0.80$ MPa(g)         & 设计压力（外压） & 全真空（FV，$-0.1013$ MPa(g)） \\
    工作压力         & $0.60$ MPa(g)         & 安全阀整定压力   & $0.85$ MPa(g) \\
    设计温度         & $170$ ℃               & 工作温度         & $155$ ℃（正常）/ $100$ ℃（温水充液） \\
    MDMT             & $-15$ ℃ @ $0.85$ MPa(g) & 全真空金属温度 & $100$ ℃ \\
    \midrule
    操作介质         & 冷却水（冷媒）        & 介质组别     & 第二组（无毒、不易燃） \\
    介质密度         & $1000$ kg/m³          & 腐蚀裕量     & $3.0$ mm（壳体/封头） \\
    有效容积         & $5.5$ m³              & 充装系数     & $\le 0.90$ \\
    \midrule
    罐体内径 $D_i$   & $1800$ mm             & 筒体长度 $L$（T/T） & $1700$ mm \\
    筒体名义厚度 $\delta_n$ & $14$ mm       & 封头名义厚度 $\delta_{nh}$ & $14$ mm \\
    封头型式         & EHA 标准椭圆封头      & 封头标准     & GB/T 25198—2010\cite{gbt25198-2010} \\
    焊接接头系数 $\phi$ & $0.85$             & 无损检测     & 局部 RT $\ge 20\%$，Ⅲ 级合格 \\
    \midrule
    壳体/封头材料    & Q345R 正火钢板        & 材料标准     & GB/T 713.1—2023 \\
    接管材料         & 20 号钢无缝钢管       & 接管标准     & GB 6479—2013 \\
    裙座材料         & Q345R                 & 地脚螺栓     & M30 $\times$ 400，$n=8$ \\
    设备净重         & $\approx 4245$ kg     & 操作总重     & $\approx 7245$ kg \\
    水压试验压力 $P_T$ & $1.02$ MPa          & 试验介质     & 洁净纯水（$Cl^- \le 25$ mg/L） \\
    \midrule
    基本风压 $q_0$   & $450$ N/m²（50 年一遇） & 地面粗糙度 & B 类 \\
    抗震设防烈度     & 7 度（$0.10g$）      & 设计地震分组 & 第一组，Ⅱ 类场地 \\
    设备布置方式     & 露天，立式，裙座支承  & 保温/保冷    & 不保温 \\
    \bottomrule
  \end{tabularx}
\end{table}
